\chapter{Downstream Utility: VaR Backtesting and QLIKE}
\label{ch:downstream}

\begin{tcolorbox}[feynbox={Sky}{BBg},title={Before and After}]
\textbf{Before this chapter you need:} GARCH and QMLE
(\S\ref{sec:garch}, \S\ref{sec:qmle}), hypothesis testing, $p$-values and
statistical power (\S\ref{sec:hypothesis}--\S\ref{sec:power}), VaR and expected
shortfall (\S\ref{sec:var}--\S\ref{sec:es}), TSTR (\S\ref{sec:tstr}), the
step-budget argument (\S\ref{sec:sgd}).\\
\textbf{After it you will understand:} why conditional VaR discriminates where
unconditional VaR cannot, what Kupiec and Christoffersen each test, why QLIKE
inverts at small samples, and a confound in this project's own walk-forward
design that made fold number and training budget inseparable.
\end{tcolorbox}

The preceding chapter asked whether synthetic data \emph{looks} right. This one
asks whether it \emph{works}: fit a risk model to synthetic returns, apply it to
real ones, and see whether the forecasts hold up.

%% ─────────────────────────────────────────────────────────────────────────────
\section{Unconditional and Conditional VaR}
\label{sec:condvar}

\situation{
Two flood warnings. One says the river exceeds 3 metres on 5\% of days --- true
every day of the year, and equally uninformative every day. The other says
tomorrow it will exceed 4.2 metres, because it has been raining all week.
}

\problem{
Both are called VaR and they test entirely different things. One of them is
useless for evaluating a generator, and knowing which --- and why --- determines
whether the downstream metric measures anything at all.
}

\workedarith{
\textbf{Unconditional VaR.} Take the $5$th percentile of the whole sample. For
BOVESPA's 20-year data that is about $-2.7\%$, and it is the same number every
day.

\medskip
\textbf{Conditional VaR.} Let GARCH supply a volatility forecast for tomorrow:
\[
  \hat\sigma_t^2 = \hat\omega + \hat\alpha r_{t-1}^2 + \hat\beta\hat\sigma_{t-1}^2 ,
  \qquad
  \mathrm{VaR}_t = \hat\mu + \hat\sigma_t z_{0.05} , \quad z_{0.05} = -1.645 .
\]
On a quiet day this gives roughly $-1.5\%$; after a large move, roughly $-3.5\%$.
It moves with the market.

\medskip
\textbf{Now the decisive comparison.} Coverage error is
$|\hat{v} - \alpha|$, where $\hat{v}$ is the observed exceedance fraction.

\medskip
\begin{tabular}{@{}lr@{}}
\toprule
Series & \texttt{var\_coverage\_error} \\
\midrule
Gaussian iid noise, matched $\sigma$ & $0.0276$ \\
Real GARCH on real data              & $0.0259$ \\
Shuffled real data                   & $0.0227$ \\
\bottomrule
\end{tabular}

\medskip
\textbf{Read those three rows.} Gaussian white noise, a correctly specified GARCH
model, and a shuffled deck all score within $0.005$ of each other. The metric
cannot separate a genuine volatility model from noise, so it discriminates
nothing.

\medskip
\textbf{Why.} Unconditional coverage depends only on the marginal distribution ---
does the $5$th percentile sit in the right place? That is a fidelity question,
already answered by Wasserstein and quantile MSE. Any series with the right
marginal passes, including one with no temporal structure whatsoever.

\medskip
\textbf{What conditional VaR adds.} It asks whether the forecast rises \emph{on
the right days}. Shuffled data has the correct marginal and no information about
which days are turbulent, so its conditional forecasts are worthless even though
its unconditional coverage is fine. The discrimination lives in the timing, which
is what Christoffersen tests in \S\ref{sec:christoffersen}.
}

\formaldef{
\textbf{Value at Risk} at level $\alpha$:
\[
  \mathrm{VaR}_\alpha = \inf\{x : \Prob(r_t \leq x) \geq \alpha\} .
\]
For GARCH with normal innovations,
$\mathrm{VaR}_{\alpha,t} = \hat\mu + \hat\sigma_t\Phi^{-1}(\alpha)$.

\textbf{Coverage error} is $|\hat{v} - \alpha|$ with $\hat{v}$ the observed
exceedance fraction. Real GARCH on real BRICS data gives $\hat{v} = 7.59\%$, hence
$|0.0759 - 0.05| = 0.0259$.
}

\analogybreak{
Even conditional VaR reports only a threshold. It says nothing about losses beyond
it --- that is expected shortfall (\S\ref{sec:es}) --- and it inherits the normality
assumption in $z_{0.05}$, which \S\ref{sec:clt} showed to be badly wrong in the
tails. Using $-1.645$ on data with $\hat\alpha \approx 3$ understates the true
quantile. The backtests below partly compensate by testing realised coverage
rather than trusting the model.
}

\whyhere{
This is why \texttt{var\_coverage\_error} sits in the descriptive group rather
than among the ranked temporal metrics: the three-row table above shows it failing
to separate a shuffled deck from a real GARCH fit. The conditional machinery,
tested by Kupiec and Christoffersen together, is what carries the downstream
signal.
}

%% ─────────────────────────────────────────────────────────────────────────────
\section{The Kupiec Backtest}
\label{sec:kupiec}

\situation{
Your VaR model promises exceedances on $5\%$ of days. Over 200 days you observe
40, which is $20\%$. Something is wrong. Over 200 days you observe 12, which is
$6\%$ --- is that wrong too, or just luck?
}

\problem{
Every observed exceedance rate differs from $5\%$ by something. A test is needed
to separate ordinary sampling variation from genuine miscalibration --- and, as
\S\ref{sec:power} warned, to know when the test can detect anything at all.
}

\workedarith{
\textbf{The statistic.} With $n$ test days, $x$ observed exceedances,
$\hat{p} = x/n$ and claimed $\alpha$:
\[
  LR_{uc} = 2\ln\bigl[\hat{p}^{x}(1-\hat{p})^{n-x}\bigr]
          - 2\ln\bigl[\alpha^{x}(1-\alpha)^{n-x}\bigr] \sim \chi^2(1) .
\]
Reject at $5\%$ when $LR_{uc} > 3.84$.

\medskip
\textbf{A worked case.} $n = 200$, $\alpha = 0.05$, so $10$ exceedances are
expected.

At $x = 12$, $\hat{p} = 0.06$:
\[
  \ln\bigl[0.06^{12}(0.94)^{188}\bigr] = 12\ln 0.06 + 188\ln 0.94 = -33.76 - 11.63 = -45.39 ,
\]
\[
  \ln\bigl[0.05^{12}(0.95)^{188}\bigr] = 12\ln 0.05 + 188\ln 0.95 = -35.95 - 9.64 = -45.59 .
\]
\[
  LR_{uc} = 2(-45.39 + 45.59) = 0.40 .
\]
Far below $3.84$: not rejected. Twelve exceedances out of 200 is unremarkable.

\medskip
\textbf{Now the power problem.} Suppose the true rate is $7\%$ while the model
claims $5\%$ --- a genuinely broken model, exceeding its VaR $40\%$ more often
than advertised. Can the test find it?

\medskip
\begin{tabular}{@{}lrl@{}}
\toprule
Test days $n$ & Power & Setting \\
\midrule
$125$     & $17\%$ & 5-year, per market \\
$496$     & $\approx 56\%$ & 20-year, per market \\
$2{,}480$ & $\approx 99\%$ & 20-year, pooled \\
\bottomrule
\end{tabular}

\medskip
\textbf{At $n = 125$, a broken model goes undetected five times in six.} A
non-rejection carries essentially no information. Only at pooled sample size does
the test become a test.
}

\formaldef{
\textbf{Kupiec (1995) unconditional coverage test:}
\[
  H_0 : p = \alpha \qquad\text{against}\qquad H_1 : p \neq \alpha ,
\]
with $LR_{uc} \sim \chi^2(1)$ under $H_0$ and rejection at $LR_{uc} > 3.84$ for a
$5\%$ level.
}

\analogybreak{
Kupiec tests the \emph{total} number of exceedances and nothing about their
timing. A model whose exceedances all fall in a single catastrophic week has
perfect Kupiec coverage and is useless --- it failed in exactly the period that
mattered. That gap is precisely what the next section closes.
}

\whyhere{
This is why downstream utility is computed on the pooled test set rather than per
market. Reporting a per-market non-rejection at $17\%$ power as evidence that a
generator produces well-calibrated VaR would be a claim the data cannot support.
}

%% ─────────────────────────────────────────────────────────────────────────────
\section{The Christoffersen Independence Test}
\label{sec:christoffersen}

\situation{
A smoke alarm that goes off the right number of times a year, but always on the
same afternoon. The annual count is correct; the alarm is worthless.
}

\problem{
Kupiec is satisfied by the right \emph{number} of exceedances however they are
arranged. Volatility clusters, so a model that ignores clustering will bunch its
exceedances --- and that is the failure mode that matters most, since it is
concentrated in crises.
}

\workedarith{
\textbf{The idea.} Let $I_t = \mathbf{1}[r_t < \mathrm{VaR}_{\alpha,t}]$. Under a
correct model, whether today is an exceedance should not depend on whether
yesterday was.

\medskip
\textbf{Count the transitions.} Over $200$ days with $10$ exceedances, suppose:

\medskip
\begin{tabular}{@{}lrrr@{}}
\toprule
& $I_t = 0$ & $I_t = 1$ & total \\
\midrule
$I_{t-1} = 0$ & $184$ & $5$ & $189$ \\
$I_{t-1} = 1$ & $5$   & $5$ & $10$ \\
\bottomrule
\end{tabular}

\medskip
\textbf{Conditional exceedance rates.}
\[
  \hat\pi_{01} = \frac{5}{189} = 0.026 ,
  \qquad
  \hat\pi_{11} = \frac{5}{10} = 0.500 .
\]

\medskip
\textbf{Read those two numbers.} After a calm day, the chance of an exceedance is
$2.6\%$. After an exceedance, it is $50\%$ --- nineteen times higher. The
exceedances are heavily clustered.

\medskip
\textbf{And yet Kupiec passes.} The overall rate is $10/199 = 5.03\%$, almost
exactly the promised $5\%$. Kupiec sees a perfectly calibrated model;
Christoffersen sees one that fails whenever it is needed.

\medskip
\textbf{Why clustering is the expected failure.} Volatility clusters
(\S\ref{sec:stylisedfacts}, fact 3), so a model using a \emph{constant} VaR
threshold will inevitably breach repeatedly during turbulent stretches and never
during calm ones. Independence of exceedances is the direct test of whether the
conditional volatility model is doing its job.
}

\formaldef{
\textbf{Christoffersen (1998) independence test.} With transition counts $n_{ij}$
for $I_{t-1} = i$, $I_t = j$:
\[
  LR_{\text{ind}} = 2\ln\!\left[
    \frac{\hat\pi_{01}^{n_{01}}(1-\hat\pi_{01})^{n_{00}}\hat\pi_{11}^{n_{11}}(1-\hat\pi_{11})^{n_{10}}}
         {\hat\pi^{n_{01}+n_{11}}(1-\hat\pi)^{n_{00}+n_{10}}}\right] \sim \chi^2(1) ,
\]
where $\hat\pi_{ij} = n_{ij}/(n_{i0}+n_{i1})$ and $\hat\pi$ is the pooled rate.

The \textbf{conditional coverage} test combines both:
\[
  LR_{cc} = LR_{uc} + LR_{\text{ind}} \sim \chi^2(2) ,
\]
testing correct rate and independent timing simultaneously.
}

\analogybreak{
The test examines only first-order dependence --- whether today depends on
yesterday. Exceedances clustered at a two- or three-day lag, with no
day-to-day link, would pass. Given that volatility autocorrelation is measurable
past lag 50 (\S\ref{sec:acf}), first-order independence is a considerably weaker
requirement than genuine independence.
}

\whyhere{
$LR_{cc}$ is a substantially harder bar than Kupiec alone, and it is the test that
distinguishes a generator that learned volatility dynamics from one that only
matched the marginal. It is the downstream counterpart of the fidelity--temporal
split in Chapter~\ref{ch:evaluation}: Kupiec is the fidelity question, and
independence is the temporal one.
}

%% ─────────────────────────────────────────────────────────────────────────────
\section{QLIKE Loss}
\label{sec:qlike}

\situation{
Two forecasters. One predicts tomorrow's return, a single number you can check
against what happened. The other predicts tomorrow's \emph{volatility} --- and
volatility is never observed directly, only inferred. How do you score the second?
}

\problem{
Volatility is latent. Any scoring rule must work from a noisy proxy such as
$r_t^2$, and a badly chosen rule will reward the wrong forecast when fed a noisy
proxy. This is the problem \emph{properness} solves.
}

\workedarith{
\textbf{The loss.}
\[
  \mathrm{QLIKE}_t = \log\hat\sigma_t^2 + \frac{r_t^2}{\hat\sigma_t^2} .
\]
Two terms pulling opposite ways: forecast too high and $\log\hat\sigma_t^2$
punishes you; too low and $r_t^2/\hat\sigma_t^2$ punishes you.

\medskip
\textbf{Check that the minimum is in the right place.} Suppose the true
conditional variance is $\sigma^2 = 4$ and, on average, $r_t^2 = 4$. Evaluate the
expected loss at several forecasts:

\medskip
\begin{tabular}{@{}lrrr@{}}
\toprule
$\hat\sigma^2$ & $\log\hat\sigma^2$ & $4/\hat\sigma^2$ & total \\
\midrule
$2$ & $0.693$ & $2.000$ & $2.693$ \\
$3$ & $1.099$ & $1.333$ & $2.432$ \\
$4$ & $1.386$ & $1.000$ & $\mathbf{2.386}$ \\
$5$ & $1.609$ & $0.800$ & $2.409$ \\
$8$ & $2.079$ & $0.500$ & $2.579$ \\
\bottomrule
\end{tabular}

\medskip
The minimum is at $\hat\sigma^2 = 4$, the true value. \checkmark That is what
makes QLIKE a proper scoring rule, and it holds even though $r_t^2$ is an
extremely noisy proxy for $\sigma_t^2$.

\medskip
\textbf{The reference.} TRTR --- train and test on real data --- gives
$\mathrm{QLIKE} \approx -8.134$ on the pooled set. A synthetic dataset whose GARCH
parameters transfer well should land close to it.

\medskip
\textbf{The small-sample inversion.} At per-market 5-year sample size
($n \approx 126$):

\medskip
\begin{tabular}{@{}lr@{}}
\toprule
Series & QLIKE \\
\midrule
Gaussian noise & $-6.888$ \\
Real data      & $-6.876$ \\
\bottomrule
\end{tabular}

\medskip
\textbf{Gaussian noise scores better than real data.} The metric has inverted.
The shuffled control's standard deviation across seeds at this $n$ is $3.354$ ---
larger than any difference between models. At this sample size QLIKE is measuring
noise.

\medskip
At pooled $n \approx 2{,}480$ it is stable and interpretable, which is why it is
computed there and nowhere else.
}

\formaldef{
\textbf{QLIKE} (Patton 2011) is
\[
  \mathrm{QLIKE}(r,\hat\sigma) = \frac{1}{T}\sum_{t=1}^{T}
    \left[\log\hat\sigma_t^2 + \frac{r_t^2}{\hat\sigma_t^2}\right] .
\]
A scoring rule is \textbf{proper} if its expectation is minimised by the true
conditional variance. Patton shows QLIKE is proper and, crucially, \textbf{robust
to noisy volatility proxies} --- using $r_t^2$ in place of the unobservable
$\sigma_t^2$ does not shift the minimiser.
}

\analogybreak{
QLIKE tests whether the estimated parameters produce good variance forecasts. It
does not test whether GARCH$(1,1)$ is the right \emph{specification} --- leverage
effects and jumps are invisible to it --- and a synthetic model with wrong
parameters but the right unconditional volatility could score well by accident.
Kupiec and Christoffersen supply the independent checks.
}

\whyhere{
QLIKE is the primary downstream ranking signal: it is proper, it directly tests
whether GARCH structure learned from synthetic data transfers to real data, and at
pooled $n \approx 2{,}480$ it has adequate power. The inversion at
$n \approx 126$ is the concrete demonstration of why pooling is not an optional
convenience.
}

%% ─────────────────────────────────────────────────────────────────────────────
\section{Walk-Forward Validation}
\label{sec:walkforward}

\situation{
Judging a weather model on its forecast for one particular July tells you about
that July. Judging it across many separate periods, each forecast made without
knowledge of what followed, tells you about the model.
}

\problem{
A single train/test split gives one number with no measure of its variability, and
it may land on an unrepresentative period --- which for MOEX across 2022 it
certainly would.
}

\workedarith{
\textbf{The protocol.} Divide the series into $K$ folds. For fold $k$, train on
everything before it and test on it. Never shuffle: fold order is time order.

\medskip
\textbf{Why not shuffled $K$-fold.} Returns have serial dependence, and windows
overlap. Shuffling puts near-duplicate windows in both train and test, which is
look-ahead leakage. The measured effect on the discriminative AUC null was
$0.506 \to 0.584$ (\S\ref{sec:aucnull}) --- a shift comparable to the effects being
measured.

\medskip
\textbf{Why the full series is used.} An earlier version ran walk-forward on the
$496$-point test split alone, producing folds of $168$ to $414$ points. At
sequence length 128 and batch size 64 that is $0$ to $2$ batches per fold, so
models were evaluated essentially untrained --- and the discriminative AUC came out
at $1.000$ in every fold, a perfect score that reflected an untrained generator
rather than a detectable difference. The protocol now concatenates train,
validation and test into roughly $4{,}953$ points per market.

\medskip
\textbf{What it buys.} Mean and standard deviation across folds, fold-specific
anomaly detection --- MOEX's invasion fold is visibly different --- and a measure of
temporal stability per architecture.
}

\formaldef{
\textbf{Rolling-origin walk-forward} (CTBench protocol, Liao et al.\ 2025):
\begin{itemize}
  \item split the full series into $K$ folds ($K = 5$ here);
  \item for fold $k$, retrain from scratch on the data preceding it and evaluate
        on it;
  \item report mean $\pm$ standard deviation across folds.
\end{itemize}
Fold order is never shuffled, and each fold's model is trained from a fresh
initialisation.
}

\analogybreak{
Retraining per fold is what makes the protocol expensive and what makes it honest
--- reusing one model across folds would leak later data into earlier evaluations.
But it introduces a subtlety that the next section shows this project got wrong:
folds differ in training-set size, and unless the budget is specified carefully,
they differ in training \emph{amount} too.
}

\whyhere{
Five folds per market give a distribution rather than a point estimate, which is
what allows a claim that one architecture beats another to be more than an
artefact of a single split.
}

%% ─────────────────────────────────────────────────────────────────────────────
\section{The Walk-Forward Budget Confound}
\label{sec:foldconfound}

\situation{
A study compares teaching methods across five schools. It later emerges that the
first school ran the course for two weeks and the last for five months. The
comparison across schools is now uninterpretable --- method and duration are
entangled.
}

\problem{
Walk-forward folds have deliberately different training-set sizes: fold 0 trains
on the least data, fold 4 on the most. Under an epoch-based budget, that means
they also receive different numbers of gradient steps --- and the two effects
cannot be separated.
}

\workedarith{
\textbf{Where the confound comes from.} An epoch is
$\lfloor n_{\text{windows}}/B \rfloor$ gradient steps (\S\ref{sec:sgd}). Fold 0
has fewer windows, so at a fixed epoch count it gets fewer steps.

\medskip
\textbf{Measured in this project.}

\medskip
\begin{tabular}{@{}lrrl@{}}
\toprule
Fold & Gradient steps & Autoencoder $R^2$ & \\
\midrule
$0$ & $110$ & $+0.116$ & below the $0.8$ threshold \\
$1$ & $240$ & $+0.981$ & healthy \\
\bottomrule
\end{tabular}

\medskip
\textbf{Fold 1 received $2.2\times$ the training of fold 0.} Its $R^2$ is not
marginally better; it is the difference between an embedding that carries
information and one that does not (\S\ref{sec:timegan}).

\medskip
\textbf{Why this is a confound and not merely noise.} Fold index is confounded
with two things at once:
\begin{itemize}
  \item \emph{training budget}, which rises with fold number;
  \item \emph{market regime}, since later folds cover later periods.
\end{itemize}
If fold 4 scores better than fold 0, there is no way to attribute the improvement.
More training and a calmer regime predict the same result.

\medskip
\textbf{The fix.} Specify the budget in gradient steps and derive the epoch count
per fold:
\[
  n_{\text{epochs}} = \left\lceil\frac{\text{steps}}{n_{\text{batches}}}\right\rceil .
\]
Every fold then receives the same number of updates regardless of its size. In the
smoke run this produced $506$ and $504$ Phase-1 steps on two folds against targets
of $500$ --- equal to within the rounding forced by whole epochs.

\medskip
\textbf{What remains confounded.} Regime still varies by fold, and that is
intrinsic to walk-forward rather than a defect. The step budget removes the
\emph{avoidable} confound, leaving the one the protocol exists to expose.
}

\formaldef{
Under an epoch-based budget with expanding training windows, gradient steps grow
with fold index, so fold effects and budget effects are not identified separately.

Under a step-based budget, gradient steps are held constant across folds and the
remaining variation is attributable to data --- regime, sample composition --- alone.
}

\analogybreak{
Equal gradient steps is not equal information: fold 4 still trains on more data,
so each step of its budget sees a richer sample. The step budget equalises
\emph{optimisation effort}, not statistical power. That is the right quantity to
hold fixed for comparing architectures, and it should not be mistaken for making
the folds equivalent.
}

\whyhere{
This confound is the fold-level counterpart of the model-level one in
\S\ref{sec:sgd}, and it has the same cause: an epoch is a data-dependent unit
masquerading as a fixed one. Both are fixed by the same change. Any walk-forward
result produced before that change conflates training budget with fold, and should
not be read as evidence about temporal stability.
}
